% !TEX root = ../algebra_02.tex

\section{Нормальные подгруппы}

\begin{Remark}{Мотивация}{}
	Нормальность подгруппы (совпадение левых и правых смежных классов) является необходимым условием для корректного введения групповой операции на множестве смежных классов, то есть для построения факторгруппы.
\end{Remark}

\begin{Proposition}{Эквивалентные определения нормальной подгруппы}{}
	Пусть $H$ --- подгруппа группы $G$. Тогда следующие условия эквивалентны:
	\begin{enumerate}
		\item $\forall g \in G, \forall h \in H \colon gh = h'g$ для некоторого $h' \in H$.
		\item $\forall g \in G, \forall h \in H \colon ghg^{-1} \in H$ (элемент, сопряженный к $h$ при помощи $g$, содержится в $H$).
		\item $\forall g \in G \colon Hg = gH$ (левые и правые смежные классы совпадают).
		\item $\forall g \in G \colon gHg^{-1} \subset H$.
		\item $\forall g \in G \colon gHg^{-1} = H$.
	\end{enumerate}
\end{Proposition}

\begin{proof}
	$1 \implies 2$:
	Если $gh = h'g$ для некоторого $h' \in H$, то $ghg^{-1} = h' \in H$.
	
	$2 \implies 4$:
	По определению $gHg^{-1} = \{ ghg^{-1} \mid h \in H \}$. Из условия $ghg^{-1} \in H$ следует $gHg^{-1} \subset H$.
	
	$4 \implies 5$:
	Имеем $gHg^{-1} \subset H$ для любого $g \in G$. Пусть $g' = g^{-1}$. Тогда $(g')^{-1}Hg' \subset H$, что означает $g^{-1}Hg \subset H \implies H \subset gHg^{-1}$. Учитывая оба включения, получаем $gHg^{-1} = H$.
	
	$5 \implies 3$:
	Если $gHg^{-1} = H$, то умножая обе части равенства справа на $g$, получаем $gH = Hg$.
	
	$3 \implies 1$:
	Для любых $g \in G$ и $h \in H$ элемент $gh \in gH$. Так как $gH = Hg$, то $gh \in Hg$, следовательно, $gh = h'g$ для некоторого $h' \in H$.
\end{proof}

\begin{Definition}{Нормальная подгруппа}{}
	Подгруппа $H < G$, удовлетворяющая любому из эквивалентных свойств выше, называется \emph{нормальной подгруппой}. Обозначается $H \triangleleft G$.
\end{Definition}

\begin{Example}{Примеры нормальных подгрупп}{}
	\begin{enumerate}
		\item Тривиальные подгруппы: $\{e\} \triangleleft G$ и $G \triangleleft G$.
		\item В любой абелевой группе все её подгруппы являются нормальными.
		\item Знакопеременная группа $A_n \triangleleft S_n$ (группа четных перестановок нормальна в симметрической группе).
		\item Ядро любого гомоморфизма групп является нормальной подгруппой.
	\end{enumerate}
\end{Example}

\begin{Lemma}{Критерий нормальности для подгруппы индекса 2}{}
	Если $(G:H) = 2$, то $H \triangleleft G$.
\end{Lemma}

\begin{proof}
	Пусть $g \in G$.
	
	Если $g \in H$, то левый и правый смежные классы совпадают с самой подгруппой:
	\[
	gH = H = Hg.
	\]
	Если $g \notin H$, то так как смежных классов всего два, а их объединение дает всю группу $G$, смежный класс $gH$ состоит в точности из всех элементов $G$, не входящих в $H$. То есть $gH = G \setminus H$. Аналогичным образом правый смежный класс $Hg = G \setminus H$.
	
	Следовательно, во всех случаях $gH = Hg$, и подгруппа $H$ нормальна.
\end{proof}

\begin{Example}{Пример подгруппы, не являющейся нормальной}{}
	Рассмотрим симметрическую группу $G = S_3$ и её подгруппу $H = \langle(1\,2)\rangle = \{(1\,2), e\}$.
	
	Пусть $h = (1\,2) \in H$ и $g = (1\,2\,3) \in G$. Вычислим сопряженный элемент:
	\[
	ghg^{-1} = (1\,2\,3)(1\,2)(1\,3\,2) = (2\,3).
	\]
	Так как $(2\,3) \notin H$, элемент не сохранился в подгруппе при сопряжении. Значит, $gHg^{-1} \not\subset H$, и $H$ не является нормальной подгруппой в $G$.
\end{Example}